\textbf{1. Xác định góc mở từ mô hình tạo ảnh}

Gọi góc mở bao trọn đoạn $AB$ là $\alpha = \widehat{AMB} = \widehat{A'MH}$. Xét tam giác vuông $\Delta A'MH$, ta có hệ thức lượng giác:
\begin{equation}
    \tan \alpha = \frac{d_1}{f} \implies \alpha = \arctan\left(\frac{d_1}{f}\right)
\end{equation}

Tương tự, với góc mở $\beta = \widehat{BMC} = \widehat{HMC'}$, xét tam giác vuông $\Delta C'MH$, ta thu được:
\begin{equation}
    \tan \beta = \frac{d_2}{f} \implies \beta = \arctan\left(\frac{d_2}{f}\right)
\end{equation}

\textbf{2. Thiết lập phương trình quỹ tích trên mặt đất}

Trên bản đồ, camera $M$ nhìn đoạn thẳng $AB$ dưới góc $\alpha$. Theo định lý quỹ tích cung chứa góc, điểm $M$ phải nằm trên đường tròn $(O_1)$ đi qua $A$ và $B$. 
Độ dài đoạn thẳng $AB = \left((x_2-x_1)^2 + (y_2-y_1)^2\right)^{1/2}$. 
Bán kính đường tròn $(O_1)$ được tính bằng hệ thức: $R_1 = \frac{AB}{2\sin\alpha}$.
Tâm $O_1(a_1, b_1)$ nằm trên đường trung trực của $AB$. Bằng hình học giải tích cơ bản, tọa độ tâm $O_1$ được xác định theo công thức:
\begin{equation}
    a_1 = \frac{x_1 + x_2}{2} + \frac{y_1 - y_2}{2\tan\alpha} \quad ; \quad b_1 = \frac{y_1 + y_2}{2} + \frac{x_2 - x_1}{2\tan\alpha}
\end{equation}
Theo Định lý Pytago, điểm $M(x, y)$ cách tâm $O_1$ một khoảng $R_1$, ta có phương trình đường tròn:
\begin{equation}
    x^2 + y^2 - 2a_1 x - 2b_1 y + c_1 = 0 \quad (\text{với } c_1 = a_1^2 + b_1^2 - R_1^2)
\end{equation}

Hoàn toàn tương tự, camera $M$ nhìn đoạn $BC$ dưới góc $\beta$ nên $M$ nằm trên đường tròn $(O_2)$ đi qua $B$ và $C$. Tọa độ tâm $O_2(a_2, b_2)$, bán kính $R_2$ và hệ số $c_2$ được tính bằng cấu trúc công thức y hệt như (3). Ta có phương trình đường tròn $(O_2)$:
\begin{equation}
    x^2 + y^2 - 2a_2 x - 2b_2 y + c_2 = 0 
\end{equation}

\textbf{3. Giải phương trình tìm tọa độ điểm $M(x,y)$}

Vì $M$ là giao điểm của hai đường tròn, tọa độ $(x, y)$ phải thỏa mãn cả (4) và (5). 
Lấy phương trình (4) trừ đi phương trình (5), các thành phần bình phương $x^2$ và $y^2$ triệt tiêu, ta nhận được phương trình đường thẳng đi qua dây cung chung (đi qua trạm $B$ và camera $M$):
\begin{equation}
    2(a_2 - a_1)x + 2(b_2 - b_1)y + (c_1 - c_2) = 0 \implies y = kx + m
\end{equation}
(với $k$ và $m$ là các hệ số dễ dàng rút ra từ phương trình).

\begin{figure}[ht!]
\centering
    \begin{tikzpicture}[scale=1.3]
        % Định nghĩa tọa độ các điểm
        \coordinate (M) at (0,0);
        \coordinate (A) at (-3,4);
        \coordinate (B) at (0,5);
        \coordinate (C) at (4,3);

        % Tâm và bán kính của Đường tròn 1 và 2
        \coordinate (O1) at (-0.8333, 2.5);
        \def\Rone{2.635} 
        \coordinate (O2) at (1.25, 2.5);
        \def\Rtwo{2.795}

        % Vẽ 2 đường tròn quỹ tích
        \draw[color=blue, thick, dashed] (O1) circle (\Rone);
        \draw[color=red, thick, dashed] (O2) circle (\Rtwo);

        % Dây cung chung (Đường thẳng đi qua B và M)
        \draw[thick, purple] (0, -1.5) -- (0, 6) node[right] {\textit{Đường thẳng qua B và M}};

        % Vẽ các tia nhìn
        \draw[thick] (M) -- (A);
        \draw[thick] (M) -- (C);
        \draw[thick, gray, dotted] (A) -- (B) -- (C);

        % Đánh dấu các Drone và M
        \fill[orange] (A) circle (2.5pt) node[above left, text=black] {$A(x_1, y_1)$};
        \fill[orange] (B) circle (2.5pt) node[above left, text=black] {$B(x_2, y_2)$};
        \fill[orange] (C) circle (2.5pt) node[above right, text=black] {$C(x_3, y_3)$};
        \fill[black] (M) circle (3pt) node[below left=0.1cm] {\textbf{Máy ảnh $M(x_M, y_M)$}};
        
    \end{tikzpicture}
\caption{Đường thẳng đi qua giao điểm B và M của hai quỹ tích đường tròn.}
\end{figure}

Thế biểu thức $y = kx + m$ vào lại phương trình đường tròn (4), ta thu được một phương trình bậc hai theo ẩn $x$ có dạng:
\begin{equation}
    Px^2 + Qx + S = 0
\end{equation}

Đến đây, thay vì giải công thức nghiệm Delta phức tạp, ta nhận thấy: Vì trạm $B(x_2, y_2)$ là giao điểm hiển nhiên của hai quỹ tích, nên $x = x_2$ chắc chắn là một nghiệm của phương trình (7).
Áp dụng \textbf{Định lý Vi-ét} cho tổng hai nghiệm, ta có:
\begin{equation}
    x_M + x_2 = -\frac{Q}{P} \implies x_M = -\frac{Q}{P} - x_2
\end{equation}
Cuối cùng, thay $x_M$ vào lại đường thẳng (6), ta tính được tung độ $y_M$:
\begin{equation}
\implies y_M = k \left(-\frac{Q}{P} - x_2\right) + m 
\end{equation}

\vspace{0.5cm}

\noindent \textbf{\large Phổ điểm}
\begin{center}
\renewcommand{\arraystretch}{1.5}
\begin{tabularx}{\textwidth}{|c|X|c|}
\hline
\textbf{Phần} & \textbf{Nội dung} & \textbf{Điểm thành phần} \\ \hline
\textbf{1} & Thiết lập đúng tỉ số lượng giác $\tan \alpha$ và $\tan \beta$. & 0.5 \\ \cline{2-3} 
 & Biểu diễn được góc $\alpha$ và $\beta$ theo $d_1, d_2, f$. & 0.5 \\ \hline
\textbf{2} & Nêu được công thức tính bán kính $R_1, R_2$ và tọa độ tâm $O_1, O_2$. & 1.0 \\ \cline{2-3} 
 & Xây dựng đúng hai phương trình đường tròn $(O_1)$ và $(O_2)$. & 0.5 \\ \hline
\textbf{3} & Trừ hai phương trình để tìm ra phương trình đường thẳng bậc nhất. & 0.5 \\ \cline{2-3} 
 & Rút thế để đưa về phương trình bậc hai $Px^2 + Qx + S = 0$. & 0.5 \\ \cline{2-3} 
 & Áp dụng Định lý Vi-ét loại bỏ nghiệm $B(x_2, y_2)$, tìm đúng tọa độ $M(x_M, y_M)$. & 0.5 \\ \hline
\multicolumn{2}{|r|}{\textbf{Tổng điểm}} & \textbf{4.0} \\ \hline
\end{tabularx}
\end{center}